%%%% RESUMO
%%
%% Apresentação concisa dos pontos relevantes de um texto, fornecendo uma visão rápida e clara do conteúdo e das conclusões do
%% trabalho.

\begin{resumoutfpr}%% Ambiente resumoutfpr
Este trabalho apresenta um estudo comparativo entre três das principais ferramentas gratuitas de integração contínua: GitHub, GitLab e Jenkins. O objetivo é analisar a complexidade e a performance de cada ferramenta, buscando identificar suas vantagens e desvantagens no contexto de desenvolvimento de software. A justificativa para este estudo está na crescente adoção de práticas DevOps e a necessidade de soluções eficientes e de baixo custo para automatizar fluxos de integração contínua. A metodologia utilizada envolveu a criação de pipelines em cada uma das plataformas, seguido de uma avaliação dos tempos de execução, facilidade de configuração, suporte a plugins e a curva de aprendizado. Os testes foram realizados em um ambiente controlado, garantindo que as variáveis de infraestrutura fossem constantes para todas as ferramentas. 

\textbf{Palavras-chave}: Entrega contínua; GitHub; GitLab; Jenkins; análise de performance; Continuous Delivery; DevOps.

 
\end{resumoutfpr}
